% !TITLE 满江红·怒发冲冠
% !AUTHOR 岳飞
% !DYNASTY 两宋
% !GENRE 词
% !ORDER 15
% !EDITION 《全宋词》，唐圭璋编，中华书局
% !EDITION_STATUS pending

\begin{poem}
怒发冲冠\textsuperscript{1}，凭栏处\textsuperscript{2}、潇潇雨歇。
抬望眼\textsuperscript{3}、仰天长啸，壮怀激烈。
三十功名尘与土\textsuperscript{4}，八千里路云和月\textsuperscript{5}。
莫等闲\textsuperscript{6}、白了少年头，空悲切。

靖康耻\textsuperscript{7}，犹未雪；臣子恨，何时灭！
驾长车、踏破贺兰山缺\textsuperscript{8}。
壮志饥餐胡虏肉，笑谈渴饮匈奴血\textsuperscript{9}。
待从头、收拾旧山河，朝天阙\textsuperscript{10}。
\end{poem}

\begin{annotations}
\notetext{1}{怒发冲冠：愤怒得头发直竖，把帽子都顶了起来。语出《史记·廉颇蔺相如列传》“怒发上冲冠”。}
\notetext{2}{凭栏处、潇潇雨歇：倚着栏杆的时候，一阵急雨刚刚停歇。潇潇，形容风雨急骤的声音；歇，停止。}
\notetext{3}{抬望眼、仰天长啸：抬起头向远处望，对着天空长声呼啸。啸，撮口发出清越的长音，古人常用它来抒发胸中的郁气。}
\notetext{4}{三十功名尘与土：三十岁上下挣下的功名，像尘土一样微不足道。这是自谦，也是不甘——功业还远远不够。}
\notetext{5}{八千里路云和月：八千里的征战路途，只有天上的云和月一路相伴。写转战之远、行军之苦。}
\notetext{6}{莫等闲：不要随随便便地虚度。等闲，轻易、随便。}
\notetext{7}{靖康耻：靖康二年（1127），金兵攻破汴京，掳走宋徽宗、宋钦宗，北宋灭亡。这是宋朝最深的国耻，也是这首词全部的恨意所由来。}
\notetext{8}{贺兰山缺：贺兰山在今宁夏与内蒙古交界处，这里借指金人占据的北方。缺，山口。岳飞一生征战并未到过贺兰山，此处是借代，取其“边塞”之意。}
\notetext{9}{壮志饥餐胡虏肉，笑谈渴饮匈奴血：饿了要吃敌人的肉，谈笑间要喝敌人的血——极言对敌人的切齿之恨。胡虏、匈奴都借指金兵。}
\notetext{10}{朝天阙：朝见皇帝，向朝廷报捷。天阙，皇宫前的楼观，代指朝廷。}
\end{annotations}

\begin{appreciation}
《满江红·怒发冲冠》是岳飞最著名的一首词，也是中国古典诗词中最慷慨激烈的爱国之作之一。需要先说清楚一桩悬案：这首词不见于宋、元两代的著录，直到明代中期才广泛流传，作者是否为岳飞，学界至今有争论，余嘉锡等学者曾提出质疑，也有人为岳飞辩护。这里仍依照旧说，把它系在岳飞名下。

词从一腔怒气起笔。“怒发冲冠，凭栏处、潇潇雨歇。”怒得头发竖起顶起帽子，是《史记》里蔺相如的故事；此时他凭栏而立，一场急雨刚刚停歇。雨歇而怒不息，一个“冲冠”先把全词的情绪顶了起来。“抬望眼、仰天长啸，壮怀激烈”，抬头远望，对着天空长声呼啸——啸是古人抒发胸中积郁的方式，而胸中的壮怀太满，只能借一声长啸吐出。

“三十功名尘与土，八千里路云和月。”三十岁上下的功名，在他眼里轻如尘土；八千里路的征战，一路只有云和月相随。这一联把功名与跋涉对举：功名不值一提，可那八千里路的云月，却是实实在在的辛苦。一个“尘与土”，一个“云和月”，一轻一重，写尽了他对轩冕的冷淡和对征战的执着。“莫等闲、白了少年头，空悲切。”不要轻易虚度，等到两鬓斑白再来悲叹就迟了。这一句从家国之恨转出一个关于时间的提醒，千百年来激励过无数后来人。

下阕的恨更直白。“靖康耻，犹未雪；臣子恨，何时灭！”靖康二年，金兵攻破汴京，掳走徽钦二帝，北宋覆亡——这是宋朝立国以来最大的国耻，也是这首词全部恨意的由来。“臣子恨，何时灭”一连两个短句，几乎是咬牙切齿。“驾长车、踏破贺兰山缺。”贺兰山在西北，岳飞一生征战并未到过那里，这里取其边塞之意，借指金人盘踞的北方。“壮志饥餐胡虏肉，笑谈渴饮匈奴血”，是把切齿之恨说得最赤裸的两句，胡虏、匈奴都指金兵；写到这里，恨意已经推到了极致，未必句句都能坐实。末句一收：“待从头、收拾旧山河，朝天阙。”只等从头收拾旧日的山河，再向朝廷报捷。

历史上的岳飞，确实也走到了离成功最近的地方。绍兴十年他率岳家军北伐以郾城与颍昌两战大破金军主力，进抵朱仙镇，距汴京只余四十五里，中原义军纷纷响应。就在这时，宋高宗与秦桧连发十二道金牌召他班师，他回身南下，随即被解除兵权，次年以“莫须有”的罪名被害于风波亭，年三十九。“待从头、收拾旧山河”那六个字，终究没能等到。

《无衣》里“与子同袍”的战士，喊的是同仇敌忾的誓约；岳飞这一首，是那种誓约在一位统帅笔下的回响。读它的时候要知道它有作者之争，却也不必因此否定它——即便抛开岳飞其人，这也是一首把家国之恨写得最坦白的词。
\end{appreciation}
